\section{Discussion and Limitations}\label{sec:discussion}
With the same model identifier and action caps, the layered harness completed 41 more of the 183 public-game levels and used 35.5\% fewer actions than the baseline. Most of the reduction was in actions used on levels that each harness left unfinished (Section~\ref{sec:result-resources}), and game overs were less than half as frequent (Section~\ref{sec:result-game-overs}). These results support using the four layers together on the public games tested. If they hold under replication, they indicate that a large part of a language-model agent's score on ARC-AGI-3 depends on what evidence and notes its harness supplies at each decision.

The comparison does not identify which changes accounted for the improvement. All four layers, and the supporting features that accompanied them, changed together. The controlled tests held the other layers fixed while comparing the level-brief package, but that package itself contains four features. Separating their effects would need repeated tests that remove individual layers.

One whole-game run per condition per game does not allow an estimate of variation across repeated runs on the same game. The two harnesses also played on different days, so changes in the model backend over time could have affected the comparison. I did not conduct additional whole-game evaluation trials beyond those reported here, or whole-game tests that removed individual layers.

Equal action caps do not imply equal inference resources. The layers increased input per action, replies rejected for invalid rule-list updates required additional model calls, and calls stopped by timeouts reported no usage. The 11.0\% reduction therefore applies to reported tokens in these runs and does not establish a reduction in total compute consumed.

Prior exposure to the games further limits generalization. Before the final evaluation, I had studied the eight-game development group, inspected baseline failures across all 25 games, played levels of two games myself, and run the controlled level-brief tests with the layer code fixed (Section~\ref{sec:brief-design}). Clearing the agent's memory between games does not remove my own exposure to them. The 17 games outside the development group are therefore not a held-out test set. Differences between the groups could also reflect game difficulty.

Interruptions affected both the completion and the resource comparisons. I retained all three runs halted by timeouts. These halts may have reduced completion, and they shortened the recorded action and token totals. If the halted runs had used every remaining action, the action reduction would have been 29.8\% rather than 35.5\% (Section~\ref{sec:result-resources}).

The checks on the records confirm what the runs recorded, not how the model would act if it ran again. The continuation checks verified that reconstructed requests matched saved requests. Action replay separately checked the recorded actions and scores. For 25 \texttt{lf52} actions, the final board, game state, and completed-level count matched, while intermediate animation frames differed. Neither procedure repeated the model's inference process.

The action budget also differed from the official leaderboard's. Applying the leaderboard's per-level budget to the saved runs changed each score by less than 0.3 points (Section~\ref{sec:result-budget}). That calculation could only remove actions, so it cannot show what either agent would have done under a different budget.

The public records allow readers to reconstruct the scores and reported calculations. The full whole-game requests and replies remain local, so the public repository alone cannot reproduce every audit based on those logs or account for unreported usage.

\subsection{Priority for confirmatory evaluation}
The next priority is a new whole-game comparison that schedules the two harnesses' runs on each game close together and randomizes or balances their order, so that the date of a run cannot explain a difference. The protocol should fix the harness, prompts, game versions, action caps, and inference settings beforehand, together with matched reset, retry, and timeout policies, the number of repetitions, the primary metric, and the analysis. Each run should begin with empty memory. The protocol should also fix the exact model revision where available; otherwise, changes on the provider's side remain a limitation. Repeated paired runs, one from each harness on the same game, would estimate variation across runs on the same game. Where access permits, games that have not informed development would test transfer better than the 17 public games outside the development group. The evaluation record should retain every scheduled run, including halts and failed calls, together with all available usage reports. I have not conducted these confirmatory runs, and the results do not include them.

\section{Conclusion}\label{sec:conclusion}
On 25 public games, with one whole-game run per condition per game, the layered Claude Opus 5.5 harness scored 85.37\% in competition-mode replay, compared with 46.20\% for the baseline. It completed 166 of 183 levels, won 17 games, and used 35.5\% fewer environment actions and 11.0\% fewer reported tokens. Completion improved in 15 games, stayed the same in nine, and declined in one. Under the rule set in advance, one of the two controlled level-brief tests received the label \emph{no detectable difference} and the other the label \emph{helps}. These findings support using the four layers together on the public games. If they hold under replication, they indicate that a large part of a language-model agent's score on ARC-AGI-3 depends on what evidence and notes its harness supplies at each decision. Establishing reliability, individual effects, and transfer needs a new replication with the two harnesses run close together, tests that remove individual layers, and evaluation on previously unused games.

\section*{Reproducibility and Artifact Availability}
The parent project archives the original baseline source at \code{38f41e9d2e9937fd5bcbfbb6fe2c6e126f597039}. The \href{https://github.com/dcw06/ARC-AGI-3-Opus-5.5-Harness/tree/17d00cfdef7944dda9fd9388efa49fccce4877aa}{code and evidence archive}, reviewed at \code{17d00cfdef7944dda9fd9388efa49fccce4877aa}, contains the harness, plans, studies, replay pages, final results, usage report, log counts, game-over counts, and run figures. The plans retain the evaluated files' hashes \citep{dai2026code,dai2026harnessexport}.

The \href{https://arcprize.org/scorecards/0f882c2e-0451-41bd-86c5-211a54785c56}{baseline replay card} and \href{https://arcprize.org/scorecards/96932ddf-ab6a-4354-94b2-9644447dae14}{layered replay card} provide per-game replay and recording links \citep{dai2026baselinefinalcard,dai2026finalmanifest}. I checked public access to both cards and the archived code on November 23, 2026. On November 24, 2026, every link cited or printed in this paper opened without signing in.

The paper's source package is in the \code{paper/} folder of the same repository, on its main branch. It includes the LaTeX source, the closing replay records, per-game and overall results, token report, log counts, game-over counts, and four whole-game result summaries. The file \code{data/final_evidence_provenance.json} lists the source hashes. The command \texttt{python3 verify\_layered.py} reconstructs scores, completion counts, subgroup comparisons, replay acceptance counts, resource reductions, and the per-level budget calculation without model or network calls. The command \texttt{python3 figures/make\_result\_figures.py} redraws Figures~\ref{fig:per-game-scores}, \ref{fig:per-game-outcomes}, \ref{fig:level-action-ratio}, and~\ref{fig:brief-tests} from the same data. These checks do not rerun model inference or measure unreported usage. Earlier interim records and the baseline and Tycho appendices retain their historical scope.

\medskip
\noindent\textbf{AI assistance.} I used Claude Code to develop the harness and experiment tooling, and AI assistance to revise the manuscript and check calculations. The evaluated model and Claude Code are Anthropic products.
